\documentclass[12pt,a4paper]{article}
\usepackage[margin=24mm]{geometry}
\usepackage{fontspec}
\setmainfont{Times New Roman}
\setmonofont{Menlo}[Scale=0.84]
\usepackage{polyglossia}
\setdefaultlanguage{russian}
\newfontfamily\cyrillicfonttt{Courier New}[Scale=0.9]
\usepackage{booktabs,tabularx,url}
\usepackage[normalem]{ulem}
\setlength{\parindent}{1.25cm}
\setlength{\parskip}{3pt}
\renewcommand{\arraystretch}{1.18}

\begin{document}

% Титульный лист
\begin{titlepage}
\begin{center}
    \LARGE Национальный исследовательский университет ИТМО\\[8mm]
    Факультет Программной Инженерии и Компьютерной техники\\[37mm]
    \huge \textbf{NoSQL}\\[2mm]
    \textbf{Лабораторная работа №\,3}\\[2mm]
    \Large Вариант №\,3824
\end{center}

\vfill

\begin{flushright}
    \large Выполнили:\\
    Кузьмин Дмитрий Анатольевич,\\
    Софьин Вячеслав Евгеньевич\\[3mm]
    Группа P3309\\[3mm]
    Преподаватели:\\
    Коновалов Арсений Антонович,\\
    Деменев Тимур Гурбанович
\end{flushright}

\vfill

\begin{center}
    Санкт-Петербург 2026
\end{center}
\end{titlepage}

\section{Задание}

\textbf{Вариант №\,3824}

Продолжить разработку приложения из лабораторных работ №\,1 и №\,2. Существующие CRUD API и ранее подключённые хранилища должны продолжать работать.

Добавить wide-column систему хранения: Apache Cassandra.

\subsection*{Дополнительные параметры}

\begin{itemize}
    \item хранимые события: история изменения основных объектов;
    \item обязательный запрос: события заданного типа за период;
    \item размер временного бакета: \sout{100\,000 событий} временной период;
    \item сценарий согласованности: сравнение ONE и ALL;
    \item эксперимент с отказом: потеря доступного кворума;
    \item исследуемый механизм: размер partition и возникновение hot partition.
\end{itemize}

\subsection*{Требования к программе}

\begin{itemize}
    \item Сохранить функциональность первых двух лабораторных и записывать назначенные события в новую систему.
    \item Развернуть кластер минимум из трёх узлов с replication factor 3.
    \item Спроектировать таблицы от запросов, обосновать partition key, clustering columns и применить назначенное бакетирование.
    \item Реализовать запись событий и их чтение с сортировкой и пагинацией, включая обязательный запрос варианта.
    \item Создать отдельные денормализованные таблицы минимум для трёх запросов; использование ALLOW FILTERING запрещено.
    \item Выполнить назначенный эксперимент с уровнями согласованности и сравнить доступность и задержку операций.
    \item Провести назначенный эксперимент с отказом узла и исследовать указанный внутренний механизм.
    \item Провести нагрузочный эксперимент, зафиксировав throughput, среднюю задержку и p95.
\end{itemize}

\clearpage
\section{Реализация}

\textbf{PostgreSQL} хранит пользователей, заказы и связи преподавателей со студентами. \textbf{CouchDB} хранит каталог товаров, категории и вложения. \textbf{Riak} хранит сессии, корзины, избранное, токены восстановления и счётчики заказов, а также кэширует аналитику товаров и первую страницу истории пользователя. \textbf{Cassandra} хранит историю изменений в пространстве ключей \texttt{store} на трёх узлах с \(\mathrm{RF}=3\).

Каждое событие записывается в три таблицы. В каждой \path{event_id} типа \texttt{timeuuid} задаёт порядок от новых событий к старым.

\begin{center}
\small
\begin{tabularx}{\textwidth}{@{}lXl@{}}
\toprule
Таблица & Ключ партиции и запрос & Бакет \\
\midrule
\path{events_by_type} & (\path{event_type}, \path{bucket}). События типа за период & Сутки \\
\path{events_by_entity} & (\path{entity_type}, \path{entity_id}, \path{bucket}). История товара или заказа & 7 суток \\
\path{events_by_user} & (\path{user_id}, \path{bucket}). Действия пользователя & 30 суток \\
\bottomrule
\end{tabularx}
\end{center}

Бакет по типу самый короткий, поскольку в одну партицию попадают однотипные события по всем объектам, поэтому она может быстрее расти. Бакет по объекту длиннее, поскольку события разделены по конкретным товарам и заказам, хотя один объект могут менять разные пользователи. Для истории пользователя выбрано самое длинное окно в расчёте на более редкие действия отдельного человека. Это сокращает число бакетов, которые нужно прочитать, чтобы собрать его историю за длительный период.

\section{Эксперименты}

В опытах с согласованностью и отказом читались одни и те же 20 заранее загруженных событий, время которых укладывается в одну минуту. При записи одного события выполнялись три вставки. Измерялось время вызова репозитория. Среднее и p95 рассчитаны по успешным операциям. p95~— задержка, в которую уложились не менее 95\,\% успешных операций.

\subsection{Сравнение ONE, QUORUM и ALL}

На трёх узлах для каждой пары «уровень, операция» выполнено 5 прогревочных и 100 измеряемых попыток. При факторе репликации 3 уровни \texttt{ONE}, \texttt{QUORUM} и \texttt{ALL} требуют ответов от одной, двух и трёх реплик соответственно. В таблице приведены результаты первого из двух повторных прогонов.

\begin{center}
\begin{tabular}{@{}llrrr@{}}
\toprule
Уровень & Операция & Успех, \% & Среднее, мс & p95, мс \\
\midrule
ONE & Чтение & 100 & 1,228 & 2,167 \\
ONE & Запись & 100 & 2,129 & 3,087 \\
QUORUM & Чтение & 100 & 3,595 & 4,304 \\
QUORUM & Запись & 100 & 6,058 & 7,262 \\
ALL & Чтение & 100 & 3,755 & 4,555 \\
ALL & Запись & 100 & 6,589 & 7,316 \\
\bottomrule
\end{tabular}
\end{center}

В этом прогоне все 600 измеряемых операций прошли успешно. При \texttt{ONE} средняя задержка и p95 были минимальными для обеих операций. При \texttt{QUORUM} обе метрики были ниже, чем при \texttt{ALL}, как для чтения, так и для записи.

\subsection{Потеря доступного кворума}

Первый замер выполнен на трёх работающих узлах. Следующие замеры выполнены после остановки \texttt{cassandra3}, а затем \texttt{cassandra2}. Перед каждой серией через \texttt{nodetool status} проверялось число узлов в состоянии \texttt{UN}. Для каждой пары «уровень, операция» выполнялось 100 измеряемых попыток после 5 прогревочных.

\begin{center}
\begin{tabular}{@{}rrrr@{}}
\toprule
Живых узлов & ONE, \% & QUORUM, \% & ALL, \% \\
\midrule
3 & 100 & 100 & 100 \\
2 & 100 & 100 & 0 \\
1 & 100 & 0 & 0 \\
\bottomrule
\end{tabular}
\end{center}

Доля успешных операций была одинаковой для чтения и записи.

При \(\mathrm{RF}=3\) кворум составляет \(\lfloor 3/2\rfloor + 1 = 2\) реплики, поэтому он потерялся при отключении второго узла. \texttt{ONE} сохранил доступность вплоть до одной живой реплики. 

\subsection{Размер партиции и горячая партиция}

В основную таблицу \path{store.events_by_type} напрямую записывалось по 10\,000 синтетических событий на сценарий. В равномерном сценарии использовались десять типов событий в бакете за 1 января 2001 года UTC. В горячем сценарии события одного типа направлялись в бакет за 2 января 2001 года UTC. В обоих случаях использовались 16 параллельных клиентов и уровень \texttt{QUORUM}. Перед каждым прогоном очищались только 11 партиций эксперимента. После записи выполнялись flush и compaction на \texttt{cassandra1}, а размер партиций измерялся в SSTable инструментом \texttt{sstablepartitions}.

\begin{center}
\small
\begin{tabular}{@{}lrrrr@{}}
\toprule
Сценарий & Партиций & Строк в наибольшей & Доля её записей & Макс., байт \\
\midrule
Равномерный & 10 & 1\,000 & 10\,\% & 81\,052 \\
Горячий & 1 & 10\,000 & 100\,\% & 810\,043 \\
\bottomrule
\end{tabular}
\end{center}

При обоих порядках выполнения все записи прошли успешно. При запуске равномерного сценария первым p95 составил 5,315~мс для равномерного и 4,638~мс для горячего. При обратном порядке получено 5,484~мс и 4,959~мс соответственно.

Горячая партиция получила весь поток записей и заняла примерно в десять раз больше места, чем самая большая партиция равномерного сценария. При данной нагрузке роста p95 в горячем сценарии не наблюдалось. 

\clearpage
\subsection{Нагрузочный эксперимент}

На трёх узлах с уровнем \texttt{QUORUM} сначала измерялось чтение, затем запись. Обе операции обращались к таблице \path{store.events_by_type}. Тестовое событие выбиралось до замера, а измеряемый \texttt{SELECT} читал его по полному первичному ключу \texttt{(event\_type, bucket, event\_id)}. Каждая операция записи выполняла один \texttt{INSERT} нового события в ту же таблицу.

Для каждой операции использовались ступени параллелизма на 1, 4, 8, 16, 32, 64, 128 и 256 клиентов. На каждой ступени выполнялись 10 прогревочных и 2\,000 измеряемых операций. Throughput рассчитан как число успешных операций, делённое на длительность измеряемой серии в секундах.

\begin{center}
\small
\renewcommand{\arraystretch}{1.08}
\begin{tabular}{@{}lrrrr@{}}
\toprule
Операция & Клиентов & Throughput, оп/с & Среднее, мс & p95, мс \\
\midrule
Чтение & 1 & 270,27 & 3,674 & 4,889 \\
Чтение & 4 & 861,98 & 4,604 & 6,328 \\
Чтение & 8 & 1\,774,70 & 4,458 & 7,043 \\
Чтение & 16 & 3\,353,14 & 4,749 & 7,088 \\
Чтение & 32 & 4\,805,24 & 6,470 & 9,618 \\
Чтение & 64 & 5\,104,55 & 11,740 & 18,319 \\
Чтение & 128 & 5\,333,90 & 19,495 & 31,779 \\
Чтение & 256 & 4\,832,49 & 21,911 & 44,503 \\
\addlinespace
Запись & 1 & 283,41 & 3,499 & 4,450 \\
Запись & 4 & 915,59 & 4,337 & 6,150 \\
Запись & 8 & 2\,244,82 & 3,506 & 5,719 \\
Запись & 16 & 4\,728,06 & 3,352 & 5,002 \\
Запись & 32 & 6\,895,56 & 4,567 & 7,201 \\
Запись & 64 & 7\,445,36 & 7,967 & 12,207 \\
Запись & 128 & 7\,669,44 & 11,418 & 20,820 \\
Запись & 256 & 5\,997,86 & 13,495 & 30,769 \\
\bottomrule
\end{tabular}
\end{center}

С 64 до 128 клиентов throughput вырос лишь на 4,5\,\% для чтения и 3,0\,\% для записи. На 256 клиентах он снизился на 9,4\,\% и 21,8\,\%, а p95 вырос на 40,0\,\% и 47,8\,\% соответственно. Соответственно видим асыщение локального Docker-стенда при его текущей конфигурации и доступных ресурсах.

\section{Выводы}

В приложение добавлено хранилище истории на Apache Cassandra. Три таблицы обеспечивают выборки по типу за период, по объекту и по пользователю. Временные бакеты ограничивают период одной партиции, а \texttt{timeuuid} обеспечивает порядок от новых событий к старым и курсорную пагинацию.

Измерения показали компромисс между задержкой и доступностью уровней согласованности. \texttt{ALL} показал более высокую задержку на трёх узлах и перестал работать после первого отключения. \texttt{QUORUM} сохранил доступность при двух репликах и потерял её при одной. На локальном стенде максимальный throughput чтения и записи достигнут при 128 клиентах. При 256 клиентах обе операции показали меньший throughput и больший p95.

\end{document}
